% Direct-authority diplomatic French for printed p.152.
\SourcePage{152}

\LayerTarget{tohoku:prop:2.5.3:proof:p2}
\noindent qui est une résolution injective de $K^p$. Par suite,
$H^q(M^{p*})=H^q(\mathbf F(C(K^p)))$, et le dernier membre s'identifie à
$R^qF(K^p)$ en vertu de prop. 2.5.2.

Reste à calculer le terme médian de la deuxième ligne (2.5.5). On a tout
d'abord
\[
H_I(M)^{p,q}
=\sum_{q'+q''=q}H^p(\mathbf F^{q''}(C^{*q'}(K)).
\]
Comme $\mathbf F^{q''}$ est un foncteur exact, on obtient pour le terme
général de la somme du second membre
$\mathbf F^{q''}(H^p(C^{*q'}(K)))
=\mathbf F^{q''}(C^{q'}(H^p(K))),$
où $C(H^p(K))$ désigne une résolution injective de $H^p(K)$ (se rappeler la
définition d'une ``résolution injective'' $C(K)$ du complexe $K$ !). On
trouve donc
$(H_{II}(H_I(M)))^{p,q}
=H^q(\mathbf F(C(H^p(K))).$
qui s'identifie à $R^qF(H^p(K))$ en vertu de proposition 2.5.2, cqfd.

\LayerTarget{tohoku:construction:composite-resolving-functor-spectral-sequence}
$\mathbf F$ étant toujours un foncteur résolvant pour le foncteur exact à
gauche $F$ de $\mathbf C$ dans $\mathbf C'$, supposons de plus qu'on se donne un foncteur
covariant $G$ de $\mathbf C'$ dans une catégorie abélienne $\mathbf C''$, et que dans $\mathbf C'$
tout objet est isomorphe à un sous-truc d'un objet injectif. Pour tout
$A\in \mathbf C$, considérons la deuxième suite spectrale d'hyperhomologie du
foncteur $G$ par rapport au complexe $\mathbf F(A)$ ; c'est un foncteur
spectral en $A$, dont le terme initial, en vertu de proposition 2.5.1, se
calcule aussitôt :
\LayerTarget{tohoku:eq:2.5.6}
\[
\tag{2.5.6}
II_2^{p,q}G(\mathbf F(A))=R^pG(R^qF(A)).
\]
C'est là un isomorphisme fonctoriel avec le terme initial de la deuxième suite
spectrale du foncteur composé $GF$.

\LayerTarget{tohoku:prop:2.5.4}
\textsc{Proposition 2.5.4.} \textit{L'isomorphisme \textup{(2.5.6)}
provient d'un isomorphisme fonctoriel de la deuxième suite spectrale
d'hyperhomologie de $G$ par rapport au complexe $\mathbf F(K)$, avec la
deuxième suite spectrale du foncteur composé $GF$.}

\LayerTarget{tohoku:prop:2.5.4:proof}
Soit $C(A)$ une résolution injective de $A$, considérons $\mathbf F(C(A))$ comme un
complexe simple dans $\mathbf C'$, et considérons les homomorphismes naturels (2.5.1),
d'oú des homomorphismes correspondants pour les deuxièmes suites spectrales
d'hyperhomologie de $G$ :
\[
IIG(F(C(A)))\to IIG(\mathbf F(C(A)))
\leftarrow IIG(\mathbf F(A)).
\]
Tout revient encore à montrer que ce sont là des isomorphismes, et pour ceci,
à montrer que les homomorphismes correspondants pour les termes initiaux sont
des isomorphismes. Or en vertu de proposition 2.5.1, les homomorphismes
$H(F(C(A))\to H(\mathbf F(C(A)))\leftarrow H(\mathbf F(A))$
sont des isomorphismes, d'où aussitôt la conclusion voulue.

\LayerTarget{tohoku:diagram:2.5:functor-square}
\SourceSideDiagram{%
$\begin{array}{ccc}
\mathbf C & \xrightarrow{\ G'\ } & \mathbf C'''\\
{F,\mathbf F}\downarrow && \downarrow{F',\mathbf F'}\\
\mathbf C' & \xrightarrow[\ G\ ]{} & \mathbf C''
\end{array}$%
}

\LayerTarget{tohoku:construction:four-functor-square:p1}
Soient $\mathbf C,\mathbf C',\mathbf C'',\mathbf C'''$ quatre catégories abéliennes, on suppose pour les trois
premières que tout objet y est isomorphe à un sous-truc d'un objet injectif.
Considérons des foncteurs covariants $F,G,F',G'$ (voir diagramme ci-contre),
on suppose donné un foncteur résolvant $\mathbf F$ pour $F$ et un foncteur
résolvant $\mathbf F'$ pour $F'$, satisfaisant aux conditions de commutation
$\mathbf F'G'=G\mathbf F$\textsuperscript{7)}\SourceFootnote{7}{Il est entendu que cet isomorphisme est
fonctoriel, et respecte l'homomorphisme cobord dans les complexes
$\mathbf F'G'(A)$ et $G\mathbf F(A)$.}. Supposons de plus que $\mathbf F$ transforme
objets injectifs en

\TohokuContinuation{153}
